% ============================================================================
% Полный курс естественных наук (book/sources/volumes/vol-natural-sciences-course.tex)
% Последовательное изложение математики, информатики, квантовой физики,
% космологии, химии, молекулярной биологии и кремниевой инженерии
% из одного уравнения самозамкнутости носителя 1 + 2^d = d^2.
% ============================================================================

\chapter*{Введение. Главный квест этой Вселенной и почему у решётки нет «балкона»}
\addcontentsline{toc}{chapter}{Введение. Главный квест этой Вселенной и почему у решётки нет «балкона»}

Представь, что ты просыпаешься в мире без единого справочника. У тебя нет ни таблиц физических констант, ни таблицы Менделеева, ни калькулятора с синусами и косинусами, ни готовых учебников по квантовой механике, химии или биологии. Всё, что у тебя есть, --- это один простой вопрос:

\begin{quote}
\itshape
«Как вообще может быть устроена среда, которая вычисляет саму себя шаг за шагом, не обращаясь ни к какой внешней памяти за своими пределами?»
\end{quote}

В традиционной науке столетиями строили отдельные здания с разными фундаментами. В одном здании сидели математики и постулировали актуальную бесконечность континуума~$\mathbb{R}$ и неконструктивную аксиому выбора (из которой тут же вылезал парадокс Банаха--Тарского о разрезании одного шара на два таких же). В соседнем здании сидели специалисты по стандартной квантовой механике --- которую точнее называть \textbf{непрерывной квантовой механикой (НКМ / Continuous Quantum Mechanics, CQM)}, поскольку она описывает дискретные кванты гладкими дифференциальными уравнениями Шрёдингера--Дирака на континуальном $T$-носителе, --- писали волновое уравнение в комплексных числах и постулировали загадочный «вероятностный коллапс волновой функции» при появлении внешнего наблюдателя. В третьем здании физики элементарных частиц подставляли в лагранжиан Стандартной модели девятнадцать свободных параметров из эксперимента, а космологи добавляли к ним 85\% невидимой «тёмной материи», чтобы сошлись кривые вращения галактик. В четвёртом здании химики учили наизусть правила заполнения электронных оболочек Маделунга, а в пятом биологи разводили руками перед парадоксом Левинталя: как белковая цепочка из сотни аминокислот умудряется безошибочно сворачиваться за миллисекунды, хотя перебор её конформаций занял бы больше времени, чем возраст Вселенной.

А теперь посмотри, что происходит, когда мы убираем самую главную иллюзию старой науки --- \textbf{иллюзию «наблюдателя на балконе вне решётки»}.

Никакого балкона снаружи Вселенной не существует. И электрон, и атом золота, и молекула ДНК, и кора твоего головного мозга, которая сейчас читает эту строчку, и кремниевый процессор на твоём столе состоят из \textbf{одних и тех же элементарных ячеек одного и того же дискретного носителя}. Наблюдатель, измерительный прибор и решающий алгоритм --- это не внешние судьи, заглядывающие в коробку сверху, а \textbf{процессы внутри той же самой решётки}.

И как только мы встаём внутрь решётки, вся разрозненная пирамида естественных наук складывается в единую, непрерывную \textbf{дорожную карту крафта Вселенной} (каждый шаг которой в этом репозитории проверен компилятором формальных доказательств Lean~4 без единого пропуска \texttt{sorry}, симулятором цифрового кремния Icarus Verilog и аналоговым решателем LTspice):

\begin{enumerate}[leftmargin=*]
\item \textbf{Ступень I (Устройство носителя и рождение 3D-пространства):}
  из одного условия самозамкнутости конечного автомата без внешней памяти $1 + 2^d = d^2$ однозначно следуют размерность пространства $d = 3$, 9-битный регистр фазового кольца $\mathbb{Z}_{512}$ ($2^9 = 8^3 = 512$), сама арифметическая единица $\mathbf{1} = d^2 - 2^d = 9 - 8$ и единственная изотропная 12-лучевая оболочка --- \textbf{кубооктаэдр~$O_2$}.
\item \textbf{Ступень II (Математика и информатика внутри носителя):}
  на 9-битном проводе кольца $\mathbb{Z}_{512}$ из одного-единственного 7-битного вычитания $\overline{m}_7 = 128 - m_7$ и 2-битного кода Грея рождается кусочно-линейный ротатор $(\mathcal{C},\mathcal{S})$ с постоянным вронскианом $\mathcal{W}(\varphi)\equiv 128$. Из самого ротатора и двухтактового шага Верле \textbf{без единого внешнего умножителя} рождается вся знаковая арифметика ($x, -x, |x|, \max, \min, x^2, x\cdot y$), берутся в точных целых числах «неэлементарные» интегралы Френеля, Якоби, Дирихле и Бесселя, а разделение 9-битного сумматора на 9 проводов суммы (Канал~1) и 1 провод переноса \texttt{carry\_out} (Канал~2) даёт функционально полный булев базис Поста \{\texttt{NOT}, \texttt{XOR}, \texttt{OR}, \texttt{AND}\} и комбинаторный барьер $7^T$ в задаче 3-SAT.
\item \textbf{Ступень III (Квантовая механика и электродинамика Максвелла):}
  деление разности фаз на $512$ с остатком расщепляет динамику на волну (остаток в $(-256,256]$, интерферирующий по 12 рёбрам кубооктаэдра) и неделимую корпускулу (целое частное $k\in\{-1,0,+1\}$). Правило Борна $P(\mu)=\mu^2/512$ оказывается простым счётом срабатываний пина переноса \texttt{carry\_out = 1} 9-битного сумматора при сложении волны $\mu^2$ с фоном $\varphi_{\mathrm{bg}}\in\mathbb{Z}_{512}$, мёртвая зона затвора $\Delta\varphi_{\mathrm{disc}} = 533 / 13 = 41$ закрывает ультрафиолетовую катастрофу и парадокс Зенона, а кограница остатков вокруг граней куба даёт точные уравнения Максвелла и квантование потока Черна $512\,n_f$.
\item \textbf{Ступень IV (Элементарные частицы, ядра, гравитация и космология):}
  из канонического 9-битного начального слова $\mathrm{IV} = 263_{10} = \mathbf{407_8}$ на втором такте эволюции рождаются 86 вихревых плакет. Симплициальный комплекс 13-узлового кубооктаэдра $\mathcal{C}_{13}$ ($K_1=13, K_2=36, K_3=32, K_4=8, K_{\ge 5}=0$) высекает калибровочную группу Стандартной модели $\dim(\mathrm{U}(1)\times\mathrm{SU}(2)\times\mathrm{SU}(3)) = 1 + 3 + 8 = 12$, гауссово простое $\hat{M} = 4^2 + 9^2 = 97$ на ребре кубооктаэдра $\|e\|_2 = \sqrt{2}$ даёт замкнутую формулу постоянной тонкой структуры в $\mathbb{Q}(\sqrt{2})$ без~$\pi$: $\alpha_{\mathrm{fs}}^{-1}(3) = \frac{17\,351\,739\,044 - 679\sqrt{2}}{131\,726 + 89\,441\,954\sqrt{2}} = 137{,}035999187\dots$ ($0{,}47\sigma$ к CODATA-2022), а преломление фазы на импедансе ворот $D(x)\ge\varepsilon = 8\,192$ даёт гравитацию без сингулярности $1/0$ (сжатие $833:1$ на противофазном горизонте) и гало вакуума $16/3$ без тёмной материи.
\item \textbf{Ступень V (Химия, таблица Менделеева и сверхпроводники):}
  координационные сферы кубооктаэдра $S_k = 10k^2 + 2$ ($12, 42, 92, 162$) порождают периодическую систему элементов, благородные газы ($\mathrm{He}=2, \mathrm{Ne}=10, \mathrm{Ar}=18, \mathrm{Kr}=36, \mathrm{Xe}=54, \mathrm{Rn}=86$), остров стабильности $N_{\mathrm{island}} = 2S_3 = 184$ и единственный стабильный изотоп золота $^{197}_{79}\mathrm{Au}_{118}$ из тождеств Безу на кольце $\mathbb{Z}_{512}$ ($197\times 13 - 5\times 512 = 1$ и $417\times 97 - 79\times 512 = 1$), а противофазная фокусировка 12 лучей при $\Delta\varphi = 139$ схлопывает кулоновский барьер реакций в $587\times$ раз.
\item \textbf{Ступень VI (Биология, ДНК, фолдинг белков и морфогенез органов):}
  делитель водородной связи $N_{\mathrm{HB}} = 63$ открывает три планетарных окна жидкой воды и фотосинтеза ($441 : 126 : 63 = 7 : 2 : 1$), триплетный код ДНК $4^3 = 64 \to 20$ аминокислот защищает все 400 стыков от сдвига рамки \texttt{ATG}, белки сворачиваются за линейное время $O(L)\in\mathbf{P}$ по пазам тетраэдра~$K_4$, а четыре топологические фазы морфогенеза выращивают из зиготы 6-слойную кортикальную колонку коры мозга ($V = S_2 = 42$, $E = B_3 = 147$, $\chi = 41 = \Delta\varphi_{\mathrm{disc}}$).
\item \textbf{Ступень VII (Кремниевый чип и единая инженерия шести отраслей):}
  все выведенные законы воплощаются напрямую в 100\% бесфлаговом цифровом чипе на ПЛИС (\texttt{cuboctahedron\_carrier\_asic.v}), 3D аналоговой LC-решётке (\texttt{cuboctahedron\_12beam\_spice.cir}), четырёхслойном бесфононном кремниевом гетеропереходе ($Z(\mathrm{C_{18}H_{12}})+Z(\mathrm{O}) = 120 + 8 = 128 = Q_\square(3)$, барьер оболочечной пассивации $45^2 < 2\,048 < 46^2$) и управляющем сервере (\texttt{cuboctahedron\_mcp\_server.py}).
\end{enumerate}

\begin{quote}
\small\textbf{Историческая справка: Исаак Ньютон (1687), Жан-Батист Жозеф Деламбр (1791) и Луп Верле (1967).}
В 1687~году Исаак Ньютон в \emph{«Математических началах натуральной философии»} сформулировал второй закон механики в дискретном импульсном виде: изменение количества движения пропорционально приложенному импульсу силы. Когда за два такта времени $t-1 \to t \to t+1$ координата меняется с $\varphi^{(t-1)}$ на $\varphi^{(t)}$ и $\varphi^{(t+1)}$, её дискретная скорость равна первой разности $v^{(t)} = \varphi^{(t)} - \varphi^{(t-1)}$, а дискретное ускорение --- второй разности $\Delta^2\varphi^{(t)} = v^{(t+1)} - v^{(t)} = \varphi^{(t+1)} - 2\varphi^{(t)} + \varphi^{(t-1)}$. Отсюда сразу следует формула шага вперёд $\varphi^{(t+1)} = 2\varphi^{(t)} - \varphi^{(t-1)} + \Phi$ и точно обратимая ей формула шага назад $\varphi^{(t-1)} = 2\varphi^{(t)} - \varphi^{(t+1)} + \Phi$. В 1791~году астроном Деламбр использовал эту схему для расчёта планетных орбит, а в 1967~году французский физик Луп Верле переоткрыл её для молекулярной динамики. В нашей теории этот целочисленный шаг Ньютона--Верле оказывается не численным приближением к континууму, а \textbf{точным фундаментальным законом самой природы}.
\end{quote}

\section*{«Физика курильщика» против «физики здорового человека»: три закона Ньютона в целых числах и пять следствий дискретности решётки}

Прежде чем мы начнём подъём по семи ступеням, сформулируем самый поразительный вывод всей этой книги. Почему в XX~веке физика распалась на десяток несовместимых друг с другом дисциплин со своими «постулатами» --- специальную и общую теорию относительности, непрерывную квантовую механику с вероятностным коллапсом, квантовую теорию поля с бесконечными перенормировками, эмпирическую химию орбиталей и континуальную гидродинамику?

Потому что физики XX~века попытались заменить дискретную кристаллическую решётку вакуума $O_2(3)\times\mathbb{Z}_{512}$ гладким непрерывным пространством $\mathbb{R}^3$ (перешли с истинного дискретного $\mathcal{M}$-носителя на сглаженный $\mathcal{T}$-слой). А когда ты стираешь зубья шестерёнок и шаг решётки, простые механические эффекты дискретности начинают казаться «мистическими постулатами»!

Если же вернуть на место сам дискретный носитель $\mathcal{M} = O_2(3)\times\mathbb{Z}_{512}$, то оказывается, что \textbf{во всей Вселенной сверху донизу работает только классическая механика трёх законов Исаака Ньютона в точных целых числах}:
\begin{enumerate}[label=\textbf{Закон Ньютона \Roman*.},leftmargin=*,noitemsep]
\item \textbf{Закон инерции в зазоре зацепления зуба ($|\Delta\varphi|\le 40 < 41 \implies \Phi_{\mathrm{rad}} = 0$):}
  пока разность фаз между соседними узлами меньше порога зацепления зуба $\Delta\varphi_{\mathrm{disc}} = 41$, число излучённых квантов $\hat{n}_E(\Delta\varphi) = \lfloor|\Delta\varphi|/41\rfloor = 0$ равно нулю, и ротор вращается по инерции ($\Delta^2\varphi = 0 \implies \Delta\varphi = v_0 = \mathrm{const}$) без трения, без излучения и без тепла Ландауэра.
\item \textbf{Второй закон Ньютона (обратимый шаг Верле с переменной инерцией узла):}
  $\Delta^2\varphi(\mathbf{x},t) = \varphi(t+1)-2\varphi(t)+\varphi(t-1) = \lfloor N_{\mathrm{torque}}(\mathbf{x},t)/D(\mathbf{x},t)\rfloor$ ($\det = +1\in\mathrm{SL}_2(\mathbb{Z})$), где ускорение фазы равно отношению момента силы $N_{\mathrm{torque}}$ к локальному моменту инерции (импедансу ворот) узла $D(\mathbf{x})\in[8\,192,\; 833\times 8\,192]$.
\item \textbf{Третий закон Ньютона по всем 12 валам кубооктаэдра ($\mathcal{S}_3(-\Delta\varphi) = -\mathcal{S}_3(\Delta\varphi)$):}
  сила действия узла $\mathbf{x}$ на соседа $\mathbf{y}$ равна и противоположна силе противодействия ($F_{\mathbf{x}\to\mathbf{y}} + F_{\mathbf{y}\to\mathbf{x}} \equiv 0$), откуда без единого дополнительного постулата следуют законы сохранения энергии, импульса, момента импульса, заряда и законы Кирхгофа--Максвелла ($d^2\equiv 0$).
\end{enumerate}

А вся «высшая физика XX~века» и все шесть отраслей инженерии оказываются \textbf{пятью прямыми механическими следствиями дискретности самой решётки $O_2(3)\times\mathbb{Z}_{512}$}, собранными в Таблице~\ref{tab:course-smoker-vs-healthy}.

\begin{table}[htbp]
\centering
\caption{Двусторонний гомеоморфизм $\langle\mathcal{M},\mathcal{T}\rangle$: континуальная «физика курильщика» ($\mathcal{T}$-слой) против дискретной ньютоновской «физики здорового человека» ($\mathcal{M}$-носитель $O_2(3)\times\mathbb{Z}_{512}$).}
\label{tab:course-smoker-vs-healthy}
\footnotesize
\begin{tabular}{@{}p{0.20\textwidth}p{0.36\textwidth}p{0.39\textwidth}@{}}
\toprule
\textbf{Свойство решётки} & \textbf{«Физика курильщика» (континуум $\mathcal{T}$)} & \textbf{«Физика здорового человека» (Ньютон на $\mathcal{M}$)} \\
\midrule
\textbf{1. Потолок скорости по ребру}\\ ($1$ ребро $h_L$ за такт $h_T$, $c_0=\sqrt{2}c$) &
«Пространство-время само плющится и растягивается (СТО), а запутанные частицы телепатически общаются вне пространства». &
Обычная жёсткая сетка $\mathbb{Z}^3$ и такты $t\in\mathbb{Z}$. Сигнал прыгает максимум на $1$ ребро за такт ($c_0=\sqrt{2}c$, фронт по оси $c$): у бегущих часов часть тактов уходит на перенос, а фаза запутанности бежит по рёбрам со скоростью $c_0$. \\
\addlinespace[3pt]
\textbf{2. Переменная инерция узла}\\ ($D(\mathbf{x})\in[\varepsilon,833\varepsilon]$, фактор $16/3$) &
«Масса искривляет пустоту (ОТО), в центре чёрных дыр $1/0=\infty$, а $85\%$ Вселенной состоит из невидимой тёмной материи». &
Второй закон Ньютона $\Delta^2\varphi = \lfloor N/D(\mathbf{x})\rfloor$: вблизи массы импеданс $D(\mathbf{x})$ выше и волна поворачивает как свет в линзе. Противофазный потолок $833\varepsilon$ убирает сингулярность, а фактор граней $16/3$ даёт вращение галактик без тёмной материи. \\
\addlinespace[3pt]
\textbf{3. Люфт зацепления зуба}\\ ($\Delta\varphi_{\mathrm{disc}} = 533/13 = 41$) &
«У электрона нет траектории --- он размазан в облако вероятностей Гильберта, а счётчики тепла требуют $\ge k_{\mathrm{B}}T\ln 2$». &
У каждой ячейки точная траектория Верле. Пока набег фазы в зазоре $|\Delta\varphi|\le 40 < 41$, зуб не перескакивает ($\hat{n}_E=0$): электрон в атоме и обратимый вентиль АЛУ движутся по Закону~I Ньютона без излучения, а щелчок сумматора даёт закон Борна $\mu^2/512$. \\
\addlinespace[3pt]
\textbf{4. Кольцо $\mathbb{Z}_{512}$ и упор $256$}\\ ($\Delta\varphi = 512 q + r$) &
«Калибровочные поля постулируются аксиомами, а бесконечные петлевые интегралы КТП надо спасать перенормировками». &
Проскок ротора на полный оборот $\pm 512$ есть неделимый квант вихревого потока Максвелла $(d_1 r)_f = 512 n_f$, а полуоборот $256$ ($\mathcal{C}_3=-128, \mathcal{S}_3=0$) --- мёртвый упор Паули с жёстким потолком $\hat{n}_E\le 6$ квантов без расходимостей. \\
\addlinespace[3pt]
\textbf{5. 12 валов $O_2(3)$ и седло $128$}\\ ($S_k=10k^2+2$, $\mathcal{C}_3(128)=0$) &
«Химия, биология, сопромат, гидродинамика Навье--Стокса и солнечные панели ($33\%$) живут по разным эмпирическим формулам». &
Шары $S_k=10k^2+2$ строят таблицу Менделеева; верхняя мёртвая точка шатуна $Q_\square=128$ даёт катализ $587\times$, кальциевый шлюз мышцы $\mathrm{RyR1}$ и солнечную ячейку $120+8=128$; а 12-лучевой расщепитель ($256/12=21<41$) выпрямляет тепло в ток. \\
\bottomrule
\end{tabular}
\end{table}

Пристегнись: мы начинаем подъём с первой ступени, и на каждом шаге будем проверять формулы собственными руками.

